% !TeX program = xelatex
% Самодостатній файл UTF-8. Компіляція: xelatex <ім’я-файлу>.tex (двічі).
\documentclass[10pt,a4paper,landscape]{article}
\usepackage[margin=15mm]{geometry}
\usepackage{fontspec}
\setmainfont{FreeSerif.otf}[BoldFont=FreeSerifBold.otf,ItalicFont=FreeSerifItalic.otf,BoldItalicFont=FreeSerifBoldItalic.otf]
\usepackage{polyglossia}
\setdefaultlanguage{ukrainian}
\usepackage{longtable,array,booktabs,xcolor,fancyhdr}
\usepackage[hidelinks]{hyperref}
\definecolor{accent}{HTML}{16656B}
\pagestyle{fancy}\fancyhf{}
\fancyhead[L]{STEM • 5 клас}\fancyhead[R]{Навчальний та календарно-тематичний план}
\fancyfoot[C]{\thepage}
\setlength{\headheight}{14pt}
\setlength{\parindent}{0pt}
\setlength{\parskip}{5pt}
\renewcommand{\arraystretch}{1.18}
\newcolumntype{P}[1]{>{\raggedright\arraybackslash}p{#1}}
\begin{document}
{\LARGE\color{accent} Навчальний план STEM, 5 клас}

{\large 1 год/тиждень • 32 годин на рік}

Заклад освіти: \rule{90mm}{0.3pt}\quad Навчальний рік: \rule{30mm}{0.3pt}

Учитель/учителька: \rule{95mm}{0.3pt}\quad Клас: \rule{25mm}{0.3pt}
\section*{Пояснювальна записка}
Авторський адаптований план укладено за наданою модельною програмою «STEM. 5--6 класи (міжгалузевий інтегрований курс)» О. В. Бутурліної та О. Є. Артєм’євої (розділ 5 класу, с. 9--20 PDF) і зошитом-конспектом «STEM-LAB. 5» (О. Бутурліна та ін., за загальною редакцією О. В. Коршунової, ВД «Освіта», 2023; файл «Stem. Full.pdf»).
Мета: формувати дослідницькі, інженерні, математичні й цифрові вміння, навички співпраці та уявлення про STEM-професії через практичні проєкти.

У наданій програмі орієнтовний розподіл становить 35 годин: 3 + 7 + 5 + 5 + 5 + 5 + 5. Програма дозволяє вчителеві змінювати розподіл годин зі збереженням очікуваних результатів (с. 8 PDF). Цей план адаптовано до 32 навчальних тижнів. Збережено вступ, п’ять змістових модулів і підсумкову STEM-практику. Послідовність модулів відповідає програмі; у зошиті розділ про сили передує розділу про Всесвіт.
\textbf{Облік часу:} І семестр — $15\times1=15$ год; ІІ семестр — $17\times1=17$ год; разом 32 години. Один рядок таблиці відповідає одному уроку. Скорочення порівняно з орієнтовними 35 годинами досягнуто об’єднанням споріднених тем і використанням готових шаблонів; дослід, проєкт і презентація збережені в кожному модулі. Усі обов’язкові роботи виконуються в межах уроків; масштаб виробу добирається відповідно до часу.

\textbf{Ресурси та безпека:} зошит, папір, картон, лінійки, фарби, ліхтарик, низьковольтні навчальні набори; за можливості — комп’ютер зі Scratch і камера. Цифрові завдання мають паперовий варіант. Електричні моделі працюють лише від батарейок; не використовувати мережеву напругу. Не спрямовувати оптичні прилади на Сонце або очі. Запуски моделей проводити у вільній зоні. Складні вирізи готує вчитель.

\textbf{Джерельні позначення:} «Зошит» у таблиці — номер сторінки наданого 64-сторінкового скану, з обкладинкою як сторінкою 1. Це опора для добору вправ, а не вимога виконати всі завдання сторінок. Частину практикумів, Scratch, семестрову діагностику й хакатон адаптовано або запропоновано за програмою; вони не подаються як дослівні завдання зошита.
\newpage
\section*{Розподіл навчального часу}
\begin{tabular}{P{125mm}rrr}
\toprule Розділ & І семестр & ІІ семестр & Разом \\\midrule
Вступ & 3 & 0 & 3 \\
Людина — людина. Я у школі & 6 & 0 & 6 \\
Людина — природа. Я у Всесвіті & 6 & 0 & 6 \\
Людина — техніка. Сила — це сила & 0 & 5 & 5 \\
Людина — образ. Промінь і світло & 0 & 4 & 4 \\
Людина — знак. Під знаком STEM & 0 & 4 & 4 \\
Хакатон, фестиваль, STEM-практика & 0 & 4 & 4 \\
\midrule Разом & 15 & 17 & 32 \\
\bottomrule\end{tabular}
\section*{Оцінювання та очікувані результати}
На першому занятті — стартова бесіда й короткі практичні завдання без окремої додаткової години. Упродовж навчання — спостереження за роботою, аналіз таблиць, ескізів, моделей, короткі пояснення учнів. Після кожного проєкту — самооцінювання та взаємний відгук. Семестрові підсумки інтегровані у тижні 15 і 32.

Спільні критерії: зрозуміла мета; коректні спостереження й розрахунки; працездатність та безпечність моделі; пояснення змін після випробування; співпраця й особистий внесок. Для формувального відгуку використовуються описи «виконує з допомогою», «виконує самостійно», «пояснює та переносить у нову ситуацію». Конкретну шкалу підсумкового оцінювання визначає заклад освіти.

Наприкінці року учень/учениця формулює просту проблему й гіпотезу, планує послідовність дій, вимірює та подає дані в таблиці або діаграмі, створює і вдосконалює модель, користується знаками та цифровими засобами, пояснює рішення і називає пов’язані професії. Портфоліо містить особистий план, результати п’яти тематичних проєктів і підсумковий прототип.
\newpage
\section*{Календарно-тематичний план. І семестр}
\small\begin{longtable}{P{13mm}P{16mm}P{47mm}P{18mm}P{88mm}P{52mm}}
\toprule Тиждень / год & Дата & Модуль і тема & Зошит & Навчальна діяльність & Очікуваний результат \\\midrule\endfirsthead
\toprule Тиждень / год & Дата & Модуль і тема & Зошит & Навчальна діяльність & Очікуваний результат \\\midrule\endhead
\bottomrule\endfoot
1 / 1 & & \textbf{Вступ}\par STEM-освіта та світ професій & 2--5 & Коло асоціацій STEM; сортування професій; стартова бесіда. & Пояснює складники STEM і наводить приклади професій. \\ \midrule
2 / 1 & & \textbf{Вступ}\par Проєкт: від задуму до плану & 6--9 & Сформулювати мету шкільного STEM-центру, розподілити ролі, скласти резюме проєкту. & Визначає мету, продукт, ресурси й етапи проєкту. \\ \midrule
3 / 1 & & \textbf{Вступ}\par Мій шлях до успіху & 6--7 & Створити особистий план розвитку; обговорити рівні можливості у виборі професії. & Визначає власну навчальну ціль і спосіб перевірки поступу. \\ \midrule
4 / 1 & & \textbf{Людина — людина. Я у школі}\par Моя школа: люди та професії & 10--12 & Мініопитування про шкільні професії; таблиця й діаграма складу класу. & Збирає знеособлені дані й пояснює діаграму. \\ \midrule
5 / 1 & & \textbf{Людина — людина. Я у школі}\par Територія школи. Масштаб і план & 13--14 & Виміряти клас; обчислити площу; виконати план у вибраному масштабі. & Переносить виміри на план і застосовує умовні позначення. \\ \midrule
6 / 1 & & \textbf{Людина — людина. Я у школі}\par Умови та засоби навчання & 15--16 & Порівняти навчальні засоби різних часів; спланувати дослідження класу. & Обґрунтовує вибір засобів навчання за заданими критеріями. \\ \midrule
7 / 1 & & \textbf{Людина — людина. Я у школі}\par Чому енергоефективність важлива & 17--19 & Порівняти споживання ламп за наданими даними; запропонувати способи економії. & Обчислює витрати енергії за зразком і пояснює вибір рішення. \\ \midrule
8 / 1 & & \textbf{Людина — людина. Я у школі}\par Енергоефективна сигнальна система & 20--21 & Скласти низьковольтне коло з батарейки, вимикача та світлодіодного модуля. & Читає просту схему, збирає і перевіряє модель за інструкцією. \\ \midrule
9 / 1 & & \textbf{Людина — людина. Я у школі}\par Школа моєї мрії: захист проєкту & 22--23 & Зібрати простий макет із підготовлених деталей; представити енергоощадне рішення. & Пов’язує план, модель і потреби користувачів; оцінює внесок команди. \\ \midrule
10 / 1 & & \textbf{Людина — природа. Я у Всесвіті}\par Об’єкти Всесвіту. Сонячна система & 36--38 & Побудувати порівняльну модель планет; обговорити масштаб і межі моделі. & Розрізняє основні об’єкти й пояснює умовність їх зображення. \\ \midrule
11 / 1 & & \textbf{Людина — природа. Я у Всесвіті}\par Зоряне небо та оптичні прилади & 39--40 & Виготовити зоряний ліхтарик за шаблоном; порівняти телескоп і бінокль. & Пояснює призначення оптичних приладів і демонструє модель. \\ \midrule
12 / 1 & & \textbf{Людина — природа. Я у Всесвіті}\par Від повітряної кулі до космічної станції & 41--42 & Створити стрічку освоєння космосу; змоделювати станцію з готових елементів. & Називає функції станції та обов’язки екіпажу. \\ \midrule
13 / 1 & & \textbf{Людина — природа. Я у Всесвіті}\par Людина у відкритому космосі & 43--44 & Порівняти захисні матеріали; спроєктувати модель захисту від холоду та ударів. & Пов’язує небезпеку з властивістю захисного матеріалу. \\ \midrule
14 / 1 & & \textbf{Людина — природа. Я у Всесвіті}\par Космічні професії. Україна космічна & 45--46 & Підготувати картку професії й інформаційний буклет за матеріалами вчителя. & Називає компетентності фахівця; зазначає джерела інформації. \\ \midrule
15 / 1 & & \textbf{Людина — природа. Я у Всесвіті}\par Я у Всесвіті: представлення результатів & 36--46 & Представити космічні моделі та буклети; виконати коротку семестрову діагностику. & Пояснює рішення з опорою на дані; визначає, що варто повторити. \\ \midrule
\end{longtable}\normalsize
\newpage
\section*{Календарно-тематичний план. ІІ семестр}
\small\begin{longtable}{P{13mm}P{16mm}P{47mm}P{18mm}P{88mm}P{52mm}}
\toprule Тиждень / год & Дата & Модуль і тема & Зошит & Навчальна діяльність & Очікуваний результат \\\midrule\endfirsthead
\toprule Тиждень / год & Дата & Модуль і тема & Зошит & Навчальна діяльність & Очікуваний результат \\\midrule\endhead
\bottomrule\endfoot
16 / 1 & & \textbf{Людина — техніка. Сила — це сила}\par Сили у природі. Історія повітроплавання & 24--25 & Дослідити дію сили на рух і форму предмета; порівняти літальні апарати. & Наводить приклади дії сил і розрізняє спостереження та припущення. \\ \midrule
17 / 1 & & \textbf{Людина — техніка. Сила — це сила}\par Сила тертя & 26--28 & Порівняти рух однакового тіла на різних поверхнях; заповнити таблицю. & Робить висновок про тертя, змінюючи одну умову досліду. \\ \midrule
18 / 1 & & \textbf{Людина — техніка. Сила — це сила}\par Прості машини та механізми & 29--31 & Дослідити важіль на лінійці; знайти прості механізми у побутових предметах. & Пояснює дію важеля на основі досліду. \\ \midrule
19 / 1 & & \textbf{Людина — техніка. Сила — це сила}\par Крила: конструювання моделі & 32--34 & Виготовити паперовий літак; висунути припущення про вплив форми крил. & Створює модель за ескізом і формулює перевірювану гіпотезу. \\ \midrule
20 / 1 & & \textbf{Людина — техніка. Сила — це сила}\par Крила: випробування й захист & 32--35 & Виміряти дальність повторних запусків; порівняти результати та представити висновок. & Фіксує результати й аргументує поліпшення моделі. \\ \midrule
21 / 1 & & \textbf{Людина — образ. Промінь і світло}\par Світло і колір & 47--49 & Дослідити спектр на готовій моделі спектроскопа; порівняти кольорові зображення. & Описує спектральний склад білого світла й результати спостереження. \\ \midrule
22 / 1 & & \textbf{Людина — образ. Промінь і світло}\par Фарби та відтінки & 50--51 & Змішати фарби; створити палітру; порівняти природні барвники. & Отримує відтінки й відрізняє змішування фарб від змішування світла. \\ \midrule
23 / 1 & & \textbf{Людина — образ. Промінь і світло}\par Поезія та фотографія: образ Сонця & 52--55 & Зіставити художній образ із явищем; створити фотографію або паперову композицію світла і тіні. & Добирає засоби зображення й пояснює роль освітлення. \\ \midrule
24 / 1 & & \textbf{Людина — образ. Промінь і світло}\par Проєкт «Намалюю тобі Сонце» & 47--55 & Оформити творчу роботу з палітрою, світлиною та поясненням світлового явища. & Поєднує наукове пояснення і художній задум; презентує результат. \\ \midrule
25 / 1 & & \textbf{Людина — знак. Під знаком STEM}\par Знаки й професії системи «людина — знак» & 56--58 & Класифікувати знаки; закодувати коротке повідомлення; обговорити професії. & Розрізняє природні й штучні знаки та пояснює правила кодування. \\ \midrule
26 / 1 & & \textbf{Людина — знак. Під знаком STEM}\par Графічні знаки. Геральдика & 57--60 & Розглянути герби; створити ескіз сімейного герба з поясненням символів. & Пояснює зв’язок символу зі змістом і читає графічні повідомлення. \\ \midrule
27 / 1 & & \textbf{Людина — знак. Під знаком STEM}\par Кодування імені. STEM-символіка & 58--61 & Закодувати ім’я за обраною абеткою; створити ескіз емблеми. & Користується узгодженим кодом і обґрунтовує символи. \\ \midrule
28 / 1 & & \textbf{Людина — знак. Під знаком STEM}\par Проєкт «Я обираю STEM» & 61--63 & Оформити емблему та презентувати її значення; провести взаємооцінювання. & Передає ідеї STEM символами й пояснює власне рішення. \\ \midrule
29 / 1 & & \textbf{Хакатон, фестиваль, STEM-практика}\par STEM-практика: професії громади & 64 & Провести очну або віртуальну зустріч із фахівцем; за відсутності — дослідити підготовлений кейс. & Ставить змістовні запитання й пов’язує професію з навчальними вміннями. \\ \midrule
30 / 1 & & \textbf{Хакатон, фестиваль, STEM-практика}\par Хакатон: задум і прототип & 62--64 & Обрати проблему класу; спланувати й створити простий прототип із доступних матеріалів. & Визначає проблему, критерії успіху, ресурси й ролі. \\ \midrule
31 / 1 & & \textbf{Хакатон, фестиваль, STEM-практика}\par Хакатон: випробування і вдосконалення & 62--64 & Випробувати прототип; зафіксувати зміни; підготувати короткий виступ. & Обґрунтовує зміни результатами випробування. \\ \midrule
32 / 1 & & \textbf{Хакатон, фестиваль, STEM-практика}\par STEM-фестиваль. Підсумкова рефлексія & 62--64 & Представити портфоліо й прототип; виконати підсумкові завдання та самооцінювання. & Демонструє поступ у дослідженні, конструюванні та співпраці. \\ \midrule
\end{longtable}\normalsize
\end{document}
